% !TeX program = xelatex
% !TeX encoding = UTF-8
% ===========================================================================
% 电赛控制类报告模板 —— 小车巡线 / 自动行驶类
% 适用题号参考：2024-H、2023-I、2022-C、2021-F、2020-C、2019-A、2018-C
% 编译：xelatex template_car.tex （建议运行两次）
% ===========================================================================
% 【模板说明】
% 本模板针对"差速/舵机小车 + 巡线/循迹 + 速度闭环 + 状态机调度"这一大类
% 控制题设计。如果你的题目涉及：
%   - 黑线/电磁线/色带循迹
%   - 编码器测速 + PID 速度闭环
%   - IMU 姿态/航向辅助
%   - 多点停车、多圈、多路段调度
%   - 声光提示、无线通信、视觉识别
% 直接套用本模板，填空即可。
% ===========================================================================
\documentclass[12pt,a4paper,fontset=fandol]{ctexart}

% ---------- 页面设置 ----------
\usepackage[top=3.5cm,bottom=2.5cm,left=2.5cm,right=2.5cm]{geometry}
\usepackage{setspace}
\setstretch{1.25}

% ---------- 数学 ----------
\usepackage{amsmath}
\usepackage{amssymb}

% ---------- 表格 ----------
\usepackage{booktabs}
\usepackage{tabularx}
\usepackage{longtable}
\usepackage{array}

% ---------- 图片 ----------
\usepackage{graphicx}
\usepackage{float}
\usepackage[font=small]{caption}
\captionsetup{labelsep=space}

% ---------- 列表 ----------
\usepackage{enumitem}
\setlist{nosep}

% ---------- 页眉页脚 ----------
\usepackage{fancyhdr}
\usepackage[hidelinks]{hyperref}
\pagestyle{fancy}
\fancyhf{}
\fancyfoot[R]{\thepage}
\renewcommand{\headrulewidth}{0pt}

% ---------- 代码 ----------
\usepackage{listings}
\lstset{
  basicstyle=\ttfamily\small,
  breaklines=true,
  frame=single,
  numbers=left,
  numberstyle=\tiny,
  showstringspaces=false,
}

% ---------- 标题 ----------
\ctexset{
  section      = {format = \large\bfseries},
  subsection   = {format = \normalsize\bfseries},
}

% ---------- 快捷命令 ----------
\newcolumntype{Y}{>{\raggedright\arraybackslash}X}
\newcommand{\code}[1]{\texttt{\small #1}}
\newcommand{\blank}{\rule{4cm}{0.4pt}}   % 封面横线

\begin{document}

% =========================== 封 面 ===========================
\begin{titlepage}
  \centering
  \vspace*{2.2cm}
  {\zihao{2}\bfseries \blank 年全国大学生电子设计竞赛\par}
  \vspace{0.5cm}
  {\zihao{2}\bfseries \blank 赛区（TI杯）\par}
  \vfill
  {\zihao{1}\bfseries \blank\par}          % ← 填写作品名称
  \vspace{0.7cm}
  {\zihao{3}\bfseries 题目编号：\blank 题\par}
  \vfill
  \begin{tabular}{rl}
    \zihao{4}参赛队编号：& \rule{7cm}{0.4pt}\\[0.55cm]
    \zihao{4}参赛队学校：& \rule{7cm}{0.4pt}\\[0.55cm]
    \zihao{4}参赛队学生：& \rule{7cm}{0.4pt}
  \end{tabular}
  \vfill
  {\zihao{4}二〇\blank 年\blank 月\par}
\end{titlepage}

% =========================== 正 文 ===========================
\setcounter{page}{1}
\begin{center}
  {\zihao{2}\bfseries \blank}             % ← 再次填写作品名称
\end{center}

\begin{abstract}
\noindent
% ===== 摘 要（≤300 字）=====
% 套路：一句话点题 → 硬件组成 → 控制算法 → 创新点 → 结果指标
针对\blank 年全国大学生电子设计竞赛\blank 题中
\blank{核心任务，如：无外部标记的黑线循迹、多点停车、多圈运行}的需求，
设计了一套以\blank{MCU 型号，如 MSPM0G3507}为主控的差速轮式小车控制系统。
系统以\blank{传感器，如 8 路灰度/CCD/电磁}获得\blank{偏差信息}，
以\blank{编码器/IMU}构建\blank{速度/姿态}闭环，
结合\blank{控制算法，如纯追踪 + 增量式 PI}实现精确循迹与调速。
任务层采用\blank{路段状态机/有限状态机}，将\blank{巡线、航向对齐、直冲跨越、重捕获、圈数计数}串联，
实现全部\blank{基本要求 + 发挥部分}的统一调度。

针对\blank{题目特殊难点，如：无标记端点、失线跨越、累计误差}，
设计了\blank{解决方案，如：多条件融合过点判据 + 位姿快照校正}。
系统实现了\blank{量化结果，如：XX 秒内完成单圈，连续 X 圈成功率 X\%}。
本文给出硬件组成、关键控制公式、软件流程和可复现测试方案。

\noindent\textbf{关键词：}\blank{MCU型号}；\blank{自动循迹}；\blank{差速控制}；\blank{控制算法}；\blank{状态机}
\end{abstract}

% =========================== 第 1 章 ===========================
\section{设计方案与工作原理}

\subsection{预期实现目标定位}
本设计面向\blank{题号}题规定的场地与任务：
\begin{itemize}
  \item 场地：\blank{场地尺寸、黑线/轨迹形状、关键点位置等}
  \item 基本要求（1）：\blank{描述}
  \item 基本要求（2）：\blank{描述}
  \item 发挥部分（1）：\blank{描述}
  \item 发挥部分（2）：\blank{描述}
\end{itemize}
设计须遵守的限制条件：
\begin{itemize}
  \item 必须使用\blank{如 TI MSPM0 系列}MCU 作为主控
  \item \blank{车体尺寸、供电方式、前进方式等限制}
  \item \blank{传感器限制，如不得安装摄像头、不得使用 GPS 等}
\end{itemize}

\subsection{技术方案分析比较}

\textbf{1. 循迹方案论证。}
方案一：单路/少路红外对管循迹。元件数量少、实现简单，但仅能判断"偏左/偏右"，
无法获得横向偏差的连续量，在弧线段难以兼顾稳定性与速度。
方案二：摄像头视觉循迹。信息量最大，可提前预瞄远处轨迹，
但题目明确禁止安装摄像头（或超出预算/计算能力）。
方案三：\blank{8 路灰度/CCD 线性传感器阵列}。
通过\blank{传感器位置的加权质心}形成连续横向偏差信号，
满足题目传感器限制，计算量小，适合\blank{MCU 型号}的实时处理能力。
\emph{方案选择：}采用方案三，兼顾\blank{精度与题目限制}。

\textbf{2. 电机控制方案论证。}
方案一：开环 PWM 调速。实现最简单，但车速随电池电压和负载变化明显，重复性差。
方案二：位置式 PID。经典算法、调参资料丰富，但对积分饱和敏感，输出限幅后易出现积分累积。
方案三：增量式 PI。直接输出 PWM 增量而非绝对量，天然抗积分饱和，适合固定周期的速度闭环。
\emph{方案选择：}采用方案三，\blank{理由：实现简单、抗饱和、适合 200 Hz 固定周期}。

\textbf{3. \blank{失线/过渡段}导航方案论证。}
方案一：仅依赖灰度传感器——在无线路段完全失效。
方案二：仅依赖编码器里程计——轮胎打滑导致累积误差。
方案三：\blank{灰度特征 + 编码器里程 + IMU 航向}多传感器融合，
以 IMU 在失线段保持航向，在捕获黑线后以灰度修正。
\emph{方案选择：}采用方案三，多信息源相互校验，提升可靠性。

\subsection{系统结构工作原理}
系统由感知、控制、执行和提示四部分组成。
\blank{灰度/CCD 传感器模块}经\blank{接口方式}采样，输出\blank{轨迹位置}；
双路编码器反馈左右轮瞬时速度；
\blank{IMU}通过\blank{如 UART 115200 bps}输出\blank{航向与角速度}；
\blank{MCU 型号}在\blank{控制周期}定时中断内完成采样、状态机、速度闭环与 PWM 刷新；
\blank{电机驱动芯片}驱动两台\blank{直流减速电机}；
\blank{OLED/LED/蜂鸣器}显示运行状态并在关键点发出声光提示。

\begin{table}[H]
\centering
\caption{系统模块组成}
\begin{tabularx}{\textwidth}{p{2.2cm}p{4cm}Y}
\toprule
模块 & 主要器件/接口 & 作用\\
\midrule
主控 & \blank{MCU 型号} & 实时控制、状态机、外设管理\\
执行 & \blank{驱动芯片 + 电机} & 左右轮独立 PWM 驱动\\
速度反馈 & \blank{编码器参数} & 轮速与累计里程\\
循迹 & \blank{传感器型号/数量} & 轨迹横向偏差与捕获状态\\
姿态 & \blank{IMU 型号，接口参数} & 航向与角速度，失线段辅助\\
人机交互 & \blank{显示/声光器件} & 状态显示、关键点提示\\
\bottomrule
\end{tabularx}
\end{table}

% \begin{figure}[H]
% \centering
% \includegraphics[width=0.85\textwidth]{figures/track_layout.png}
% \caption{场地示意图}
% \end{figure}

\subsection{功能指标实现方法}
\begin{enumerate}
  \item \textbf{弧线/直线巡线：}弧线段采用\blank{纯追踪/曲率控制}模型，将横向偏差转换为角速度命令，
  并根据\blank{偏差大小/黑线覆盖数量}自动降速；直线段保持较高基础速度。
  \item \textbf{关键点检测与停车：}通过\blank{丢线特征、航向转角、里程累计和位姿邻近}多条件协同确认到达关键点，
  到达目标点后执行\blank{停车/声光提示/状态切换}。
  \item \textbf{失线段导航：}利用\blank{IMU 航向角与角速度}构造\blank{PD 航向保持}，
  设置\blank{最大距离和超时保护}，检测到黑线后连续确认重捕获再切回巡线状态。
  \item \textbf{多圈/多路段调度：}以\blank{路段状态机}统一管理巡线、对齐、跨越、停车状态，
  到达已知关键点后调用\blank{位姿快照}修正累计误差。
\end{enumerate}

\subsection{测量、控制与分析处理}
每个\blank{控制周期}控制周期中：
\begin{enumerate}
  \item 读取编码器脉冲 → 换算左右轮实时速度与累计里程；
  \item 读取\blank{循迹传感器}状态 → 计算\blank{轨迹位置}与\blank{黑线覆盖数量}；
  \item 依据当前任务状态生成线速度 $v$ 和角速度 $\omega$ 目标；
  \item 差速运动学逆解得到左右轮目标速度；
  \item 两路增量式 PI 分别计算左右 PWM 占空比并刷新输出。
\end{enumerate}
\blank{IMU}数据以\blank{更新周期}更新，用于\blank{航向对齐、直冲段保持和位姿积分}。

% =========================== 第 2 章 ===========================
\section{核心部件与电路设计}

\subsection{关键器件性能分析}
\begin{itemize}
  \item \textbf{\blank{主控 MCU 型号}：}\blank{内核+主频}，提供\blank{定时器/PWM/GPIO 中断/UART}等资源，
  满足\blank{控制频率}控制与多外设并行需求。
  \item \textbf{\blank{电机驱动芯片}：}\blank{双 H 桥}，可独立控制左右电机方向与 PWM 占空比，
  \blank{持续/峰值电流}满足电机需求。
  \item \textbf{编码器：}\blank{线数}霍尔编码器，配合\blank{减速比}减速比，
  按\blank{倍频方式}计数时每轮一周对应\blank{计数值}个计数；
  以轮周长 $C = \blank{}$\,m 和\blank{控制频率}周期可换算实时车速。
  \item \textbf{\blank{循迹传感器}：}\blank{检测原理、分辨率、采样方式}。
  \item \textbf{\blank{IMU 型号}：}\blank{输出数据类型、帧格式、更新率}。
\end{itemize}

\subsection{电路结构工作机理}
电机方向引脚控制\blank{驱动芯片 IN1/IN2}，定时器输出\blank{PWM 频率与范围}。
\blank{循迹传感器}通过\blank{地址复用/独立 I/O}读取各通道检测值。
编码器 A/B 相\blank{上升沿/双边沿}触发 GPIO 中断计数。
\blank{IMU}采用\blank{帧格式（如 0x55 帧头+校验）}解析姿态数据。
电机电源与 MCU 逻辑电源应\blank{共地/隔离}，并复核\blank{电流能力/抗干扰}。

\subsection{运动学与控制模型}

\subsubsection*{差速运动学}
设轮距为 $B$、车体线速度为 $v$、角速度为 $\omega$，则左右轮速度目标为：
\begin{equation}
  v_L = v - \frac{\omega B}{2}, \qquad v_R = v + \frac{\omega B}{2}
  \label{eq:kinematic}
\end{equation}

\subsubsection*{循迹偏差提取}
\blank{传感器}的加权位置为：
\begin{equation}
  y = \frac{\sum(s_i \cdot p_i)}{\sum s_i}
  \label{eq:weighted}
\end{equation}
式中 $s_i$ 为第 $i$ 路检测值（\blank{0/1 二值化或灰度值}），$p_i$ 为该传感器横向坐标。

\subsubsection*{纯追踪曲率控制}
将横向偏差 $y$ 转换为弧线曲率 $\kappa$ 和角速度指令：
\begin{equation}
  \kappa = \frac{2y}{L^2 + y^2}, \qquad
  \omega = -K_{\text{steer}} \cdot v \cdot \kappa
  \label{eq:pure_pursuit}
\end{equation}
式中 $L$ 为前视距离，$K_{\text{steer}}$ 为转向增益。

\subsubsection*{增量式 PI 速度控制}
以 $e(k) = v_{\text{target}}(k) - v_{\text{measured}}(k)$ 为速度偏差：
\begin{equation}
  \text{PWM}(k) = \text{PWM}(k-1) + K_p\bigl[e(k) - e(k-1)\bigr] + K_i e(k)
  \label{eq:pi}
\end{equation}
输出限幅 $\pm \text{PWM}_{\max}$，防止控制量超出驱动器有效范围。

\subsubsection*{IMU 航向保持（PD）}
失线段使用 PD 控制保持目标航向 $\psi_{\text{target}}$：
\begin{equation}
  \omega = K_{p\psi}(\psi_{\text{target}} - \psi) + K_{d\psi} \dot{\psi}
  \label{eq:heading_pd}
\end{equation}

\subsection{关键参数与接口}
\begin{table}[H]
\centering
\caption{控制系统关键参数}
\begin{tabularx}{\textwidth}{p{3.2cm}p{3.6cm}Y}
\toprule
参数 & 取值 & 说明\\
\midrule
控制周期 & \blank{如 5\,ms} & 定时中断，控制频率 \blank{如 200\,Hz}\\
轮距 $B$ & \blank{m} & 差速运动学参数\\
轮周长 $C$ & \blank{m} & 编码器速度换算\\
编码器计数 & \blank{/rev} & \blank{线数×倍频×减速比}\\
基础/弯道速度 & \blank{mm/s} & 分档速度策略\\
前视距离 $L$ & \blank{m} & 纯追踪参数\\
$K_p$/$K_i$ & \blank{} & 增量式 PI 参数\\
PWM 频率与限幅 & \blank{如 10\,kHz，$\pm$7800} & 定时器配置\\
\blank{IMU}串口 & \blank{如 115200\,bps} & 姿态数据输入\\
\bottomrule
\end{tabularx}
\end{table}

% =========================== 第 3 章 ===========================
\section{系统软件设计分析}

\subsection{系统总体工作流程}
上电后完成\blank{外设初始化（GPIO、Timer、UART、PWM、中断优先级）}和\blank{参数加载}。
系统采用\blank{"定时中断 + 主循环"}架构：
\begin{itemize}
  \item \textbf{\blank{定时中断（控制频率）}：}
  编码器测速 → \blank{IMU}数据更新 → 任务状态机 →
  \blank{循迹/导航}计算 → 目标轮速 → PI 输出 → PWM 刷新。
  \item \textbf{主循环：}
  \blank{OLED}刷新、串口协议处理、按键扫描、\blank{遥测/DataScope}数据上报。
\end{itemize}

若\blank{传感器（如 IMU）}离线，系统应\blank{降级处理策略：如暂停依赖该传感器的功能、提示故障}，
不应将无效数据视作可靠的导航依据。

\subsection{任务状态机设计}
题目规定\blank{场地特殊约束，如：不得设置额外标记、A/B/C/D 无识别特征}，
因此\blank{关键点}不能通过\blank{十字黑线/特殊标记}直接识别。
软件将\blank{任务名称}分解为以下路段状态：

\begin{table}[H]
\centering
\caption{任务状态机}
\small
\begin{tabularx}{\textwidth}{p{2.6cm}p{2.2cm}Yp{2.5cm}}
\toprule
状态 & 输入依据 & 主要动作 & 退出条件\\
\midrule
\code{\blank{巡线状态1}} & \blank{轨迹位置、位姿} & \blank{沿轨迹巡线} & \blank{丢线且接近目标点}\\
\code{\blank{对齐状态}}   & \blank{航向、角速度} & \blank{对齐目标航向} & \blank{航向误差连续满足阈值}\\
\code{\blank{跨越状态}}   & \blank{航向、重捕获} & \blank{按目标航向前进} & \blank{重捕获轨迹线}\\
\code{\blank{巡线状态2}} & \blank{轨迹位置、位姿} & \blank{沿轨迹巡线} & \blank{丢线且接近目标点}\\
\code{STOP}               & \blank{圈数/里程计数}  & \blank{速度清零、声光提示} & \blank{任务完成}\\
\bottomrule
\end{tabularx}
\end{table}

\subsection{关键模块程序设计}
\begin{itemize}
  \item \code{\blank{Tracking\_GetPos()}}：\blank{循迹传感器采样 → 加权质心计算轨迹位置 → 偏差输出}
  \item \code{\blank{Encoder\_GetSpeed()}}：\blank{脉冲计数 → 速度换算 → 累计里程}
  \item \code{\blank{IMU\_GetData()}}：\blank{帧头校验 → 解析欧拉角/角速度 → 超时检测}
  \item \code{\blank{PurePursuit\_Steer()}}：\blank{偏差 y → 曲率 κ → 角速度 ω}
  \item \code{\blank{Incremental\_PI()}}：\blank{速度偏差 → PWM 增量 → 限幅输出}
  \item \code{\blank{Task\_StateMachine()}}：\blank{状态转移条件判断 → 动作执行 → 圈数/里程更新}
\end{itemize}

\subsection{可靠性处理}
\begin{enumerate}
  \item \textbf{传感器抗噪：}\blank{如：灰度连续多次一致才判定丢线/重捕获；IMU 帧头校验+校验和}
  \item \textbf{状态切换防抖：}\blank{如：航向对齐需连续 N 个周期满足阈值才切换}
  \item \textbf{超时/距离保护：}\blank{如：失线段设置最大直冲距离 + 对齐超时，避免永久失线}
  \item \textbf{累计误差抑制：}\blank{如：到达已知关键点后执行位姿快照，将 IMU 积分结果修正为场地坐标}
  \item \textbf{输出安全限幅：}\blank{如：PWM 限幅 0～PWMmax，速度目标饱和限制}
\end{enumerate}

% =========================== 第 4 章 ===========================
\section{竞赛工作环境条件}

\subsection{软件环境}
\begin{itemize}
  \item IDE / 编译工具：\blank{如 Keil uVision5}
  \item 目标工程：\code{\blank{工程名}}
  \item 外设配置工具：\blank{如 TI SysConfig}
  \item 调试与可视化：\blank{如 串口助手、VOFA+、DataScope}
\end{itemize}

\subsection{硬件平台与场地}
硬件平台包括：
\begin{itemize}
  \item \blank{MCU 控制板型号}
  \item \blank{电机驱动芯片 + 直流减速电机 ×2}
  \item \blank{编码器参数} ×2
  \item \blank{循迹传感器型号/配置}
  \item \blank{IMU 型号}
  \item \blank{显示模块（OLED/LCD）、LED、蜂鸣器}
  \item \blank{供电方案（电池类型、容量、电压）}
\end{itemize}
测试场地：\blank{场地尺寸、轨迹形状、线宽、颜色、有无额外标记等}。

\subsection{调试条件与安全事项}
\begin{enumerate}
  \item 架空/空载状态确认：左右轮方向、PWM 极性、编码器计数符号和急停逻辑；
  \item 低速调试：完成\blank{循迹传感器}阈值标定、传感器安装高度/位置确认及控制增益整定；
  \item 电机电源应满足峰值电流需求；调试中避免车轮被障碍物卡死；
  \item 改动供电或接线后，重新验证共地与极性；上电前检查电池电压。
\end{enumerate}

% =========================== 第 5 章 ===========================
\section{作品成效总结分析}

\subsection{测试方案}
测试流程按由简到难设计：
\begin{enumerate}
  \item \textbf{单元测试：}编码器测速准确性、PWM 输出线性度、\blank{循迹传感器}响应一致性；
  \item \textbf{单段测试：}\blank{如单弧线/单直线巡线，确认轨迹跟踪精度与速度稳定性}；
  \item \textbf{单圈测试：}\blank{完整一圈，检查路段衔接、航向对齐和声光提示}；
  \item \textbf{多圈/极限测试：}\blank{连续多圈/规定路线，统计总用时、成功率、停车误差}。
\end{enumerate}
每组条件至少重复 \textbf{3 次}，记录电池电压和场地状态。
测试仪器：\blank{秒表（具体型号）、卷尺（量程/精度）、示波器（Tektronix MDO3024）、万用表（FLUKE 17B+）}。

\begin{table}[H]
\centering
\caption{实车测试记录表}
\small
\begin{tabularx}{\textwidth}{p{2.3cm}Yp{1.15cm}p{1.15cm}p{1.15cm}p{1.35cm}}
\toprule
测试项目 & 评价指标 & 第1次 & 第2次 & 第3次 & 结论\\
\midrule
基本要求（1）& \blank{指标} & 待实测 & 待实测 & 待实测 & 待判定\\
基本要求（2）& \blank{指标} & 待实测 & 待实测 & 待实测 & 待判定\\
发挥部分（1）& \blank{指标} & 待实测 & 待实测 & 待实测 & 待判定\\
发挥部分（2）& \blank{指标} & 待实测 & 待实测 & 待实测 & 待判定\\
\bottomrule
\end{tabularx}
\end{table}

\subsection{成效与不足分析}
\subsubsection*{系统成效}
\blank{逐项对照题目要求，说明达标情况。用实测数据支撑，如：
"由表 X 可见，系统单圈用时均值 X 秒，连续 X 圈成功率 X\%，满足题目要求。"}

\subsubsection*{存在不足}
\blank{实事求是分析问题，如：
（1）灰度传感器对光照变化敏感，强光/阴影下阈值需重新标定；
（2）编码器里程计在急加速/打滑时存在累积误差；
（3）IMU 长时间运行存在航向漂移，需定期在已知点校正。}

\subsection{创新特色与展望}
\subsubsection*{创新特色}
\begin{enumerate}
  \item \blank{"循线优先 + 惯导补偿 + 已知点校正"的三层导航架构}
  \item \blank{纯追踪曲率控制 + 偏差自适应分档调速}
  \item \blank{多条件协同过点判据（丢线 + 航向 + 里程 + 位姿邻近）}
\end{enumerate}

\subsubsection*{改进方向}
\begin{enumerate}
  \item \blank{自适应速度规划：根据路段曲率在线调整速度曲线}
  \item \blank{传感器自动标定：灰度阈值/IMU 零偏在线校准}
  \item \blank{黑匣子遥测回放：赛后分析误差分布，迭代控制参数}
\end{enumerate}

% =========================== 第 6 章 ===========================
\section{参考资料及文献}
\begin{enumerate}[label={[\arabic*]}]
  \item \blank{年份}年全国大学生电子设计竞赛\blank{题号}题试题。
  \item \blank{MCU 型号技术参考手册}，Texas Instruments / STMicroelectronics。
  \item \blank{电机驱动芯片数据手册}。
  \item \blank{IMU 使用说明书}。
  \item \blank{项目内部文档引用，如：docs\_202X/X题路线分解.md}。
\end{enumerate}

% =========================== 第 7 章 ===========================
\section{附件材料}

\subsection{工程文件与关键源代码}
工程文件位于 \code{\blank{firmware/工程名/}}。

核心源代码文件：
\begin{itemize}
  \item \code{\blank{control.c/h}} — 控制中断、巡线算法、速度闭环、状态机
  \item \code{\blank{motor.c/h}}   — 电机 PWM 驱动
  \item \code{\blank{encoder.c/h}} — 编码器脉冲计数与速度换算
  \item \code{\blank{sensor.c/h}}  — 循迹传感器驱动与偏差计算
  \item \code{\blank{imu.c/h}}     — IMU 数据解析与航向输出
  \item \code{\blank{board.c/h}}   — 板级初始化和引脚映射
\end{itemize}

\subsection{核心源代码清单（节选）}
% \begin{lstlisting}[language=C, caption={控制中断主函数}]
% // 在此粘贴核心控制代码
% \end{lstlisting}

\subsection{最终提交前检查清单}
\begin{itemize}
  \item[$\square$] \textbf{格式核验：}封面和正文无学校、代码、姓名；页面上方空白 $\ge$ 3cm；A4 纸 6--8 页。
  \item[$\square$] \textbf{内容核验：}测试仪器均已标注具体型号；每个方案均有 $\ge$ 2 个对比方案。
  \item[$\square$] \textbf{数据核验：}所有实测数据已补录，无虚构。
  \item[$\square$] \textbf{规则核验：}MCU 型号、传感器、供电、车体尺寸符合题目限制；无禁用器件。
  \item[$\square$] \textbf{完整性核验：}所有图、表、公式有编号；参考文献完整；页码连续。
\end{itemize}

\end{document}
